% Applications section for "The Unchosen Adjacency" (revised)
% Assumes booktabs, array, amsmath, and the definitions file.

\section{Applying the Formalism}\label{sec:apps}

The definitions are worth having only if they discriminate. A framework that condemned every interface would be a grievance in notation. This section applies the defective-action test to four cases and reports where it returns a verdict, where it withholds one, and where the case falls outside the framework.

\subsection{The verdict rule}

An action $a\in\mathcal{A}$ is judged defective only if three findings are made.
\begin{enumerate}
  \item \textbf{Consequence.} $M_a$ contains a consequential element: material attribution or association among a sufficient share of qualified observers, loss or aversion above the de minimis level, or a shortfall in an intended change that the user rates above that level.
  \item \textbf{Feasibility.} $S_a$ is cleanly realizable at $\mathcal{A}^{*}$. A comparator system that already offers such an action is the strongest evidence.
  \item \textbf{Unavailability.} $S_a$ is not cleanly realizable at $\mathcal{A}$. This is a claim about documented controls and should be dated.
\end{enumerate}
A case that fails the first finding is not a mismatch problem. A case that fails the second is a limitation of the technology and not of the interface. A case that fails the third is already served. The third finding requires the first for every available action, so evidence of consequence is the shared bottleneck. A verdict is \emph{undetermined} when a required finding has not been made.

\subsection{Four cases}

\begin{table}[h]
\centering
\small
\renewcommand{\arraystretch}{1.3}
\begin{tabular}{@{}>{\raggedright\arraybackslash}p{2.3cm}>{\raggedright\arraybackslash}p{2.6cm}>{\raggedright\arraybackslash}p{3.6cm}>{\raggedright\arraybackslash}p{2.0cm}>{\raggedright\arraybackslash}p{3.3cm}@{}}
\toprule
Case & Intended $S_a$ & Mismatch & Claim type $\kappa$ & Provisional verdict \\
\midrule
Block or profile controls to stop comments
  & Bar one person's comments on public posts; state the norm; set a duration
  & Block removes their view of public activity ($U^{-}$). Silent tools give no notice, reason, or duration ($D$)
  & Interpersonal rejection
  & Provisionally defective on personal profiles: feasibility shown on other Meta surfaces; unavailability rests on the documentation reviewed \\
Ad or recommended content beside or after a post or Reel
  & Publish the post or Reel
  & Advertisement placed among the author's posts, overlaid on a Reel, or shown immediately after it, and platform-selected content that follows it in sequence ($U^{+}$)
  & Commercial endorsement or association
  & Undetermined: feasibility supported, unavailability provisionally supported, consequence unmeasured \\
Presence as support
  & Publish the post
  & Use of the post for recommendation or growth beyond the chosen audience ($U^{+}$)
  & Institutional approval
  & Consequential for some user types by aversion. Defect status undetermined: feasibility and unavailability unverified \\
Reels delivered to a viewer
  & None (no user action)
  & None
  & Not applicable
  & Outside the framework \\
\bottomrule
\end{tabular}
\caption{Verdicts under the defective-action test. All verdicts are provisional on the evidence noted in the text.}
\end{table}

\paragraph{Blocking and the controls around it.} The intended change is a restriction on one person's participation, and in the correction case it includes stating the norm and an end to the restriction. A block removes more than participation: the surplus is a loss of perception, and the user may value the perception that a block removes, since the point of removing a comment is often to keep the surrounding conduct visible. Comparators show that a clean action is technically feasible, including within Meta's own products. Discord's timeout lets a member read while barring them from sending messages, adding reactions, or joining voice or video, records a reason, and can be lifted early \cite{discord-timeout}. Facebook Groups let administrators temporarily suspend a member from posting, commenting, and reacting \cite{fb-groups}. Facebook Pages can ban a person from posting, liking, and commenting while the public Page remains visible to them \cite{page-ban}. Instagram's Restrict hides a person's comments from everyone but the commenter unless the account owner approves them \cite{ig-restrict}. These are moderator, administrator, or owner powers on different surfaces, so they establish feasibility and not a like-for-like feature of personal profiles.

For personal Facebook profiles the controls found are of a different kind. Meta's Help Center describes the Restricted list as limiting what a person sees, namely public information and posts in which they are tagged, and not private posts \cite{fb-restricted}. That limits perception, the opposite of the separation sought. Users can also limit who may comment on public posts to Public, Friends, or those mentioned \cite{fb-comments}, which sets an audience and not a person. Several secondary guides describe a Facebook Restrict that hides a restricted person's comments. That behavior is documented for Instagram's Restrict, and the guides appear to conflate the two, so the essay does not rely on them. On the documentation reviewed, no per-person action bars commenting on a personal profile while preserving view, so the third finding is supported for that intention, provisionally, since Meta's Help Center could not be searched exhaustively.

The shortfall issue applies to the tools that exist elsewhere. Instagram's Restrict is silent: the restricted person is not told, and no reason is stated \cite{ig-restrict}. No duration is described either, and removal is manual, so it is an indefinite reversible state and not a time-bounded suspension with automatic restoration. A silent, open-ended separation conceals and does not correct, so notice, explanation, and a bounded duration are intended changes that it does not make, and they belong to $D_a$. The verdict for personal profiles is therefore a defective bundle, provisionally, with a possible defective shortfall in any silent replacement.

\paragraph{Advertising adjacency.} Advertisements can be adjacent to a user's content in three documented ways. Advertiser-facing sources describe a Facebook profile feed placement that shows advertisements in between the posts of a person's profile \cite{nestscale}, and Meta's brand-safety tools let advertisers keep their advertisements off specific profiles \cite{emarketer}. On Facebook Reels, overlay banner ads are layered over the video, at the bottom of the Reel, and post-loop ads are video ads shown between Reels \cite{smt-reels}. The overlay shares the frame of the user's content, and the post-loop advertisement follows it immediately, so adjacency here is spatial, in-frame, and temporal. An overlay layered on the video is displayed on the screen, so a screenshot of the Reel captures it, and the adjacency travels with any copy shared outside the platform. A screenshot on file shows the case: a banner labeled ``Ad'' over the lower part of a Reel, above the creator's name, audio credit, and follow control. The banner carries a control that lets the viewer dismiss it, and an ordinary author has no documented equivalent.

Sequence adjacency is not confined to advertisements. A Reel is followed in the viewing sequence by whatever the recommendation system selects next, and Meta's Reels stream is described as algorithmically recommended \cite{smt-reels}. The author has not chosen what follows, and where the following item comes from a source that neither the author nor the viewer follows, neither has chosen the connection. This adds a surplus of the same form, association with platform-selected content, and it can be consequential by aversion for authors who object to what tends to follow, whether or not any viewer attributes the following content to them. A user may decline to publish Reels for this reason. Reports of that kind are the stated reasons that the third prediction of the withdrawal section would collect, and they are not evidence for the hypothesis by themselves. Whether an author can publish a Reel that is not followed by platform-selected content was not found and is not shown here.

Whether such adjacency is consequential depends on how viewers read it, and two questions should be kept apart. The first is attribution: whether viewers take the author to have chosen or endorsed the advertisement, which is a claim about $J_o$. The second is association: whether viewers experience the advertisement as linked to the content even while knowing the author did not choose it, which is a claim about $A_o$. The likely finding is that attribution of choice is low and association is not, and if so the second is where the consequence lies. Neither has been measured. An industry-cited survey reports that a majority of respondents associate a brand with the content around its advertisements on social media \cite{peer39}. That concerns the association of a brand with content and not of an author with an advertisement, so it is suggestive of transfer through adjacency and does not measure the quantity at issue. The advertising industry's own rationale for adjacency controls is that association is real and matters, and that rationale applies to authors as well as brands.

Feasibility is supported from the advertiser side. Meta documents inventory filters with Expanded, Moderate, and Limited settings, publisher and content block lists, and third-party content block lists for Feed and Reels placements \cite{meta-brand}. Some creators have a control of their own: a 2022 report describes an opt-out of banner ads in Reels through Creator Studio for creators in Meta's in-stream monetization program, with placement of sticker ads left to the creator and selection of the advertisement left to Facebook \cite{sej-overlay}. That report is dated and secondary, but it indicates that an adjacency control for content owners has been built for monetized creators. For ordinary users, Meta's Help Center states that ads cannot be blocked entirely \cite{fb-block-ads}, and that ad preferences change which ads a person sees and not how many \cite{fb-ad-prefs}. The same page lists a setting that adjusts who can see a person's social interactions alongside ads, an existing control over one association between a user and advertising. No control that lets an ordinary profile owner shape the advertising placed beside their own posts or Reels was found. The search was not exhaustive, and other available actions, such as publishing to a restricted audience, have not been checked for whether they avoid adjacency. The asymmetry the essay can defend is narrow and stated as an absence of documented control: advertisers can decline particular profiles and some monetized creators can opt out of overlays, while no symmetrical control for ordinary posters was found. The verdict is undetermined: consequence is unmeasured, and unavailability rests on absence in the documentation reviewed.

\paragraph{Presence as support.} Posting is itself a contribution of content and attention to the platform, and that contribution is partly constitutive of using it, so it is not obviously an effect bundled with publication. The narrower candidate surplus is the use of the post for recommendation, promotion, growth, or training beyond the audience the author selected. Two separations should be distinguished here. One keeps a post from being distributed as a recommendation to people who do not already follow the author. The other keeps interactions with the post from serving as signals in recommending other things. A platform may support one without the other, and the effect set should say which is meant. Observer attribution is unlikely to establish consequence here: the claim type is institutional approval, where the threshold is plausibly high and observers may not read participation as endorsement. Aversion can establish it for some users. A user who objects to sustaining a platform they reject bears an added relation whether or not anyone reads it as endorsement, and for users whose aversion exceeds the de minimis level the surplus is consequential. Consequence for some user types does not settle a defect. A clean action, such as publishing to a chosen audience without the material entering recommendation, may exist. Meta already sets eligibility for recommendation by category of content under its Recommendation Guidelines \cite{fb-recs}, so a content-level eligibility gate exists at platform level, but a user-set opt-out for one's own posts was not found, and its availability is not shown here. The framework therefore declines to condemn the case, which shows that the test does not prove too much.

\paragraph{Reels delivered to a viewer.} A recommendation stream delivers content to the viewer. No action of the viewer is paired with a mismatch, so there is nothing for $S_a$, $U_a$, and $D_a$ to describe. It is an imposition on perception and is excluded by the definitions. The framework thereby separates mismatch in a user's own actions from unsolicited importation, a topic treated in the author's essays on proxy capture in platform interfaces, ``Safe Mode'' and ``Affect Before Architecture''. The exclusion applies to the viewer. Publishing a Reel is a user action, and the overlay and post-loop advertisements that accompany it, together with the platform-selected content that follows it, belong to the adjacency case above. Interaction with a delivered item, such as watching, can be modeled as an action with a training-signal surplus, but that analysis is not pursued here.

\subsection{A condition for the ordering}

The thresholds and proportions are normative, so it is worth stating exactly what the ordering of two verdicts depends on.

\begin{proposition}[Monotone ordering]
Suppose that, over a common population of qualified observers, the attribution distribution for advertising adjacency first-order stochastically dominates that for presence as support, so that for every $t$ the share of observers with $J_o\ge t$ is at least as large for adjacency. If $\theta_{\text{comm}}\le\theta_{\text{inst}}$ and $\rho_{\text{comm}}\le\rho_{\text{inst}}$, then whenever presence as support is consequential by misattribution, advertising adjacency is as well.
\end{proposition}

\begin{proof}
Let $F_{\text{ad}}(t)$ and $F_{\text{pr}}(t)$ be the shares with $J_o\ge t$. Then
$F_{\text{ad}}(\theta_{\text{comm}})\ge F_{\text{ad}}(\theta_{\text{inst}})\ge F_{\text{pr}}(\theta_{\text{inst}})\ge\rho_{\text{inst}}\ge\rho_{\text{comm}}$,
using that $F$ is non-increasing, dominance, the assumption that presence is consequential, and the ordering of the proportions.
\end{proof}

The proposition rests on two assumptions: that adjacency is at least as attributable as mere presence in a common observer population, and that institutional approval requires at least as much attribution and agreement as commercial endorsement. Under them, no choice of parameters reverses the ordering. It does not show that the two cases receive different verdicts, since both may be consequential or neither. Whether the assumptions hold, and how large the attribution gap is, are empirical and normative questions that the essay does not settle.

\subsection{What remains to be established}

The verdicts depend on findings not made in this essay: primary-source confirmation of the scope of Facebook's controls on personal profiles, since Meta's Help Center was only partly retrievable; viewer-side confirmation that advertisements are placed among ordinary users' posts and after their Reels (an overlay on a Reel is shown by a screenshot on file); and a small screenshot study of how viewers read them, separating four readings (the advertisement was selected by the author, endorsed by the author, merely associated with the author, or selected entirely by the platform), which are different claims about $J_o$ and $A_o$; whether any available action, such as restricted-audience publishing, already avoids adjacency; whether a user-level opt-out of recommendation use exists or is feasible; whether an author can publish a Reel that is not followed by platform-selected content; and the choice of parameters $\theta_\kappa$, $\rho_\kappa$, $\lambda_t$, and $\lambda^{+}_t$. The table is offered as a structure for those findings and not as their result.
